%%%%%%%%%%%%%%%%%%%%%%%%%%%%%%%%%%%%%%%%%%%%%%%%%%%%%%%%%%%%%%%%%%%%%%
%                                                                    %
% On a theorem of Erd\H{o}s concerning additive functions,           %
% by                                                                 %
% Jos\'e Luis Ansorena and Juan Luis Varona                          %
%                                                                    %
% Crux Mathematicorum 33 (2007), no. 5 (September), 294--296.        %
%                                                                    %
% DATE: November 20, 2006                                            %
%                                                                    %
% Written with LaTeX2e with amsart classes                           %
%                                                                    %
%%%%%%%%%%%%%%%%%%%%%%%%%%%%%%%%%%%%%%%%%%%%%%%%%%%%%%%%%%%%%%%%%%%%%%


\documentclass[12pt,a4paper]{amsart}

% \usepackage{amsthm,amsfonts}

%% PAPER WRITTEN WITH CONCRETE FONTS:
\usepackage[T1]{fontenc}
\usepackage{ccfonts,eulervm}

\newcommand{\R}{\mathbb{R}}
\newcommand{\N}{\mathbb{N}}

\newtheorem{theorem}{Theorem}
\newtheorem*{lemma}{Lemma}
\theoremstyle{remark}
\newtheorem*{rem}{Remark}

%%%%%%%%%%%%%%%%%%%%%%%%%%%%%%%%%%%%%%%%%%%%%%%%%%%%%%
\title{On a theorem of Erd\H{o}s concerning additive functions}
%%%%%%%%%%%%%%%%%%%%%%%%%%%%%%%%%%%%%%%%%%%%%%%%%%%%%%

\author[J. L. Ansorena]{Jos\'e Luis Ansorena}
\address{Departamento de Matem\'aticas y Computaci\'on,
Universidad de La Rioja,
Calle Luis de Ulloa s/n,
26004 Logro\~no, Spain}
\email{joseluis.ansorena@dmc.unirioja.es}

\author[J. L. Varona]{Juan Luis Varona}
\address{Departamento de Matem\'aticas y Computaci\'on,
Universidad de La Rioja,
Calle Luis de Ulloa s/n,
26004 Logro\~no, Spain}
\email{jvarona@dmc.unirioja.es}
\urladdr{http://www.unirioja.es/dptos/dmc/jvarona/welcome.html}

\subjclass[2000]{Primary 11A25}
\keywords{Additive functions, completely additive functions, Erd\H{o}s}

%\date{November 17, 2006}
% \newcommand{\datename}{\textit{Date}:} %% <- Original
\renewcommand{\datename}{THIS PAPER HAS BEEN PUBLISHED IN:}
\date{\textit{Crux Mathematicorum} \textbf{33} (2007), no.~5 (September), 294--296.}


%%%%%%%%%%%%%%%%%%%%%%%%%%%%%%%%%%%%%%%%%%%%%%%%%%%%%%
\begin{document}
%%%%%%%%%%%%%%%%%%%%%%%%%%%%%%%%%%%%%%%%%%%%%%%%%%%%%%

\begin{abstract}
Erd\H{o}s proved that every increasing additive
function must be a constant multiple of the logarithmic function.
We prove a weaker result that assumes that the function is 
completely additive.
In particular, what this paper does show is how wide the gulf is
between additive and completely additive functions:
proving the result for completely additive functions is very easy,
but Erd\H{o}s's proof for merely additive functions
required a formidable effort.
\end{abstract}

%%%%%%%%%%%%%%%%%%%%%%%%%%%%%%%%%%%%%%%%%%%%%%%%%%%%%%
\maketitle
%%%%%%%%%%%%%%%%%%%%%%%%%%%%%%%%%%%%%%%%%%%%%%%%%%%%%%

An \textit{additive} function is defined as a real-valued
arithmetic function $f \colon \N \to \R$ such that 
$f(nm) = f(n)+f(m)$ for all pairs of coprime integers $n$ and~$m$.
If $f(nm) = f(n)+f(m)$ for all $n,m \in \N$, then
we say that $f$ is a \textit{completely additive} function.

In~\cite[Theorem~XI, p.~17]{Erdos}, Erd\H{o}s states that
if $f$ is an additive function such that $f(n+1) \ge f(n)$
for all $n \in \N$, then $f(n) = C\log n$ for a constant $C \in \R$.
Without the hypothesis $f(n+1) \ge f(n)$, this is not true in general, as is
shown, for instance, by the completely additive function $\Omega(n)$ defined by
$\Omega(p_1^{a_1} p_2^{a_2} \cdots p_k^{a_k}) = a_1 + a_2 + \cdots + a_k$,
where $n = p_1^{a_1} p_2^{a_2} \cdots p_k^{a_k}$ is the prime decomposition
of~$n$.
Erd\H{o}s's Theorem is a deep and interesting result, but its proof is 
rather complicated; see~\cite[p.~133]{Ruzsa} and 
\cite[\S\,8.33 and 8.34, p.~265 and ff.]{Erdos-Suranyi} 
for adequate comments and additional references, including a proof 
(due to Moser and Lambek) 
that simplifies the original proof of Erd\H{o}s. 
%	In Erd\H{o}s's paper, the theorem is deduced as a consequence
%	of a deeper result, whose proof is not easy (a.k.a.,
%	\cite[Theorem~V, p.~3]{Erdos}). 

In this note, we pose a weaker result, but with a very elementary proof.
Also, we will show a nice consequence.

\begin{theorem}
\label{th:cadd}
Let $f$ be a completely additive function such that $f(n+1) \ge f(n)$
for all $n \in \N$. Then, there is a real constant $C \ge 0$
such that $f(n) = C\log n$ for all $n \in \N$.
\end{theorem}

\begin{proof}
We claim that, if $f$ and $g$ are two functions
satisfying the conditions, they must satisfy
\begin{equation}
\label{eq:fgfg}
  f(n)g(2) = f(2)g(n)\quad \forall n \in \N.
\end{equation} 
In particular, this would be true for $g(n) = \log n$.
Then, our claim implies that $f(n) = C \log n$ with $C = f(2)/\log 2$
and so the theorem is proved. Thus, we only need to check~(\ref{eq:fgfg}).

Let $n \in \N$ and $k\in\N$. Let us take $l\in \N$ such that
$2^{l-1} \le n^k < 2^l$. Then $(l-1)f(2) \le k f(n) \le lf(2)$.
The same inequality is true for the function~$g$; we can write it as
$-lg(2) \le -k g(n) \le -(l-1)g(2)$.
Multiplying the first expression by $g(2)$, the second by $f(2)$,
adding them, and dividing by~$k$, we get
$$
  -\frac{1}{k}\,f(2)g(2) \le f(n)g(2) - f(2)g(n) \le \frac{1}{k}\,f(2)g(2).
$$
Since this happens for every $k\in\N$, equation~(\ref{eq:fgfg}) follows
and the proof is complete.
\end{proof}

As a consequence, let us establish the following result:

\begin{theorem}
\label{th:NN0}
Let $f \colon \N \to \N \cup \{0\}$ be an increasing and completely additive
function. Then, $f$ is the zero function.
\end{theorem}

\begin{proof}
Let us suppose that $f$ is not the zero function.
By Theorem~\ref{th:cadd}, there is a real constant $C>0$ 
such that $f(n) = C\log n$ for all $n \in \N$.
If $f$ takes integer values, then
$f(n)/f(m) = (\log n) / (\log m)$ for all integers $n, m \ge 2$.
Therefore, $(\log n) / (\log m) = a/b$ with $a, b \in \N$.
This implies that $n^b = m^a$. But this is impossible if $n$ and $m$ 
have a non common prime factor. 
% For instance, if $n$ and $m$ are different prime numbers.
\end{proof}

\begin{rem}
By using Erd\H{o}s's Theorem, we see that the same proof serves to 
establish the corresponding result for additive functions.
\end{rem}

For completeness, let us show that, without using Theorem~\ref{th:cadd}, another
proof of the Theorem~\ref{th:NN0} can be given. Let us begin with

\begin{lemma}
Let $f$ be an increasing and completely additive function.
Then, $f(1) = 0$. Moreover, if $f$ is not the zero function, it satisfies
$f(n) > 0$ for all $n > 1$, and $f$ is strictly increasing.
\end{lemma}

\begin{proof}
We have $f(n) = f(1 \cdot n) = f(1) + f(n)$ and so $f(1) = 0$.
If $f$ is not the zero function, there exists $a \in \N$ 
such that $f(a) \ne 0$.  
Since $a \ge 1$, it follows that $f(a) \ge f(1) = 0$; hence, $f(a)>0$.
Now, given $n > 1$, there exists $k$ such that $n^k > a$; then
$k f(n) = f(n^k) \ge f(a) > 0$, and we have $f(n) > 0$.

Finally, let us suppose that $f(n) = f(m)$ with $n < m$.
This is not possible if $n = 1$, because $f(1) = 0$ and $f(m)>0$;
thus, we may assume that $1 < n < m$.
Let us take $k$ large enough such that $n^{k+1} < m^k$
(it suffices to take $k > (\log n) /(\log m - \log n)$).
Then
$$
  f(n^k) = kf(n) = kf(m) = f(m^k).
$$ 
Consequently, $f(r) = f(n^k) = f(m^k)$ for every $r$ such that 
$n^k \le r \le m^k$. In particular, $f(n^{k+1}) = f(n^k)$, and so
$(k+1)f(n) = kf(n)$, which is impossible, because $f(n) > 0$.
\end{proof}

Using this lemma we get the following.

\begin{proof}[Second proof of Theorem~\ref{th:NN0}]
For every $n \in \N$, there exists an integer $k$ such that 
$2^{k+1} - 2^k = 2^k > n$.
Let us take $n$ intermediate numbers $r_i$ between $2^k$ 
and $2^{k+1}$; that is,
$$
  2^k < r_1 < r_2 < \cdots < r_n < 2^{k+1}.
$$
Now, let us suppose that $f$ is not the zero function.
By the lemma, $f$ is strictly increasing which implies that
{%% LOCAL CHANGE. TRICK TO AVOID amsmath AND align* ENVIRONMENT
\setlength\arraycolsep{0pt}
\begin{eqnarray*}
  k f(2) &\null=\null& f(2^k) < f(r_1) < f(r_2) < \cdots
  < f(r_n) < f(2^{k+1}) \\
  &\null=\null& (k+1) f(2) = k f(2) + f(2).
\end{eqnarray*}
}% %% THE "%" IS IMPORTANT
Then, by the pigeonhole principle, we have $f(2) > n$.
% ALTERNATIVE:
% Then, there are, at least, $n$ intermediate numbers between
% $k f(2)$ and $k f(2) + f(2)$, and so $f(2) > n$.
But this happens for every $n \in \N$, which is absurd.
\end{proof}


%%%%%%%%%%%%%%%%%%%%%%%%%%%%%%%%%%%%%%%%%%%%%%%%%%%%%%
\begin{thebibliography}{9}
%%%%%%%%%%%%%%%%%%%%%%%%%%%%%%%%%%%%%%%%%%%%%%%%%%%%%%

\bibitem{Erdos}
{P. Erd\H{o}s},
On the distribution of additive functions,
\textit{Ann. of Math.} \textbf{47} (1946), 1--20.

\bibitem{Erdos-Suranyi}
{P. Erd\H{o}s and J. Sur\'anyi},
\textit{Topics in the Theory of Numbers},
Springer, 2003.

\bibitem{Ruzsa}
{I. Z. Ruzsa},
Erd\H{o}s and the Integers,
\textit{J. Number Theory} \textbf{79} (1999), 115--163.

%%%%%%%%%%%%%%%%%%%%%%%%%%%%%%%%%%%%%%%%%%%%%%%%%%%%%%
\end{thebibliography}
%%%%%%%%%%%%%%%%%%%%%%%%%%%%%%%%%%%%%%%%%%%%%%%%%%%%%%

%%%%%%%%%%%%%%%%%%%%%%%%%%%%%%%%%%%%%%%%%%%%%%%%%%%%%%
\end{document}
%%%%%%%%%%%%%%%%%%%%%%%%%%%%%%%%%%%%%%%%%%%%%%%%%%%%%%

